
\maketitle

\begin{abstract}
A consolidated reference covering everything explored in this
research programme: what I understand, what is known (and by whom),
what I proved that turned out to be classical, what remains
genuinely open, and where the understanding might lead.  Not a
paper.  A map for finding the next problem.
\end{abstract}

\tableofcontents

% ==================================================================
\section{The central question}
% ==================================================================

\emph{When does observationally coherent data determine a genuine
probability measure, and what additional commitments does this
require?}

This decomposes into three layers:
\begin{enumerate}
  \item \textbf{Extension:} compatible finitely additive charges
    on a directed system of Boolean algebras
    $\to$ $\sigma$-additive probability measure.
  \item \textbf{Realization:} measure on the dual space (Stone
    space) $\to$ measure on an underlying state space $\Omega$.
  \item \textbf{Value-definiteness:} probabilistic description
    $\to$ simultaneous definite values for all observables.
\end{enumerate}

Each layer requires a commitment the previous does not.
The algebraic structure of the observation algebra constrains
which commitments are available; distributivity is sufficient to
make all three available.

% ==================================================================
\section{Layer 1: From charges to measures}
\label{sec:extension}
% ==================================================================

\subsection{The classical results (all known)}

\begin{itemize}
  \item \textbf{Continuity from above:} A finitely additive charge
    on an algebra is $\sigma$-additive iff
    $E_n \downarrow \varnothing \Rightarrow \ell(E_n) \to 0$
    \cite[\S13, Theorem~D]{Halmos1950}.

  \item \textbf{Carath\'{e}odory extension:} A $\sigma$-additive
    premeasure on a semiring extends uniquely to a measure on the
    generated $\sigma$-algebra
    \cite[\S13]{Halmos1950}, \cite[Vol.~1, Ch.~1]{Bogachev2007}.

  \item \textbf{Kolmogorov extension:} Compatible
    finite-dimensional distributions on a product space extend to a
    unique probability measure, under standard Borel / tightness
    conditions \cite[Ch.~36]{Billingsley1995}.

  \item \textbf{Choksi inverse limits:} Compatible measures on an
    inverse system of compact Hausdorff spaces extend to the
    inverse limit \cite[Theorem~1]{Choksi1958}.

  \item \textbf{Yosida--Hewitt decomposition:} Every bounded
    finitely additive charge decomposes uniquely as
    $\ell = \ell_c + \ell_p$ where $\ell_c$ is $\sigma$-additive
    and $\ell_p$ is purely finitely additive
    \cite{YosidaHewitt1952}.
\end{itemize}

\subsection{What I proved (turned out classical)}

\textbf{Collective exhaustion (CE)} in a directed system: defined
as $\exists j \geq i$ such that $\ell_j(\pi_{ij}^{-1}(E_n)) \to 0$.
Under compatibility, the existential witness is vacuous:
$\ell_j(\pi_{ij}^{-1}(E_n)) = \ell_i(E_n)$ for all $j \geq i$.
So CE $=$ $\sigma$-additivity of each $\ell_i$.  This is not a
theorem; it is the observation that the definition contains a
redundant quantifier.

\textbf{The directed-system extension theorem}
(Paper~I, Theorem~2.3) is Kolmogorov
extension in directed-system language with ``sequential common
refinements'' replacing the usual product structure.

\textbf{The Lean formalization} (0~sorry on the Carath\'{e}odory
route, 1~sorry for Yosida--Hewitt on the Stone route) is genuine
verified infrastructure, even though the theorems are classical.

\subsection{What I understand from this}

\begin{enumerate}
  \item The passage from finite additivity to $\sigma$-additivity
    is the critical non-finitary step in probability.
    Compatibility, normalization, and finite additivity are all
    first-order / finitary conditions.

  \item No first-order condition can force $\sigma$-additivity.
    Immediate from \L{}o\'{s}'s theorem: the ultraproduct of
    Dirac masses on $\mathbb{N}$ is a purely finitely additive
    charge \cite{Hoover1978, FaginHalpernMegiddo1990}.

  \item The finite-cofinite charge on $\mathbb{N}$ is the
    canonical obstruction: normalized, finitely additive,
    not $\sigma$-additive, and arises as the ultraproduct witness.

  \item The Yosida--Hewitt decomposition gives the exact
    characterization: CE $\Leftrightarrow$ $\ell_p = 0$.
\end{enumerate}

% ==================================================================
\section{Layer 2: From dual space to realization}
\label{sec:realization}
% ==================================================================

\subsection{The classical results (all known)}

\begin{itemize}
  \item \textbf{Stone representation:} Every Boolean algebra is
    isomorphic to the clopen algebra of a compact, Hausdorff,
    totally disconnected space (its Stone space)
    \cite{Johnstone1982}.

  \item \textbf{Charge $\leftrightarrow$ measure on Stone space:}
    Every finitely additive charge on a Boolean algebra $B$
    corresponds to a regular Borel (Baire) measure on its Stone
    space $\St(B)$ (the space of ultrafilters on $B$)
    \cite[\S326]{Fremlin2003}, \cite[\S53]{Halmos1950}.

  \item \textbf{$\sigma$-additivity $\leftrightarrow$ support:}
    A charge on $B$ is $\sigma$-additive iff the corresponding
    Stone-space measure vanishes on meagre Baire sets, equivalently
    assigns full mass to countably complete ultrafilters
    \cite[Theorem~10.5.3]{RaoRao1983}.

  \item \textbf{Principal ultrafilters:} Every principal ultrafilter
    is countably complete.  The converse holds for $\mathcal{P}(\kappa)$
    when $\kappa$ is below the first measurable cardinal (Ulam's
    theorem).  For non-$\sigma$-complete algebras (typical cylinder
    algebras), countably-complete non-principal ultrafilters can
    exist without large cardinals.
\end{itemize}

\subsection{What I proved (turned out classical)}

\textbf{Stone measure exists unconditionally:} Any compatible
charges produce a Baire probability measure $\hat{\mu}$ on
$\St(\mathcal{C})$.  This is Stone representation + Choksi.
Lean-verified (0~sorry).

\textbf{Descent requires $\hat{\mu}(\pure(\Omega)) = 1$}
(where $\pure(\omega) = \{E \in \mathcal{C} : \omega \in E\}$ is
the principal ultrafilter at $\omega$)\textbf{:} concentration on
principal ultrafilters.  Under discriminability
and the assumption that all countably-complete ultrafilters in the
relevant subalgebra are principal (which holds for $\sigma$-complete
algebras below the first measurable cardinal), this is equivalent
to $\sigma$-additivity.  The support characterization is Rao--Rao
10.5.3; the principal-vs-countably-complete gap is the subtlety.

\textbf{Realization is unconstrained:} Any surjective-projecting
$S \subseteq \St(\mathcal{C})$ can serve as $\Omega$.  This is
the abstractness of Boolean algebras recognized since Stone (1936).

\subsection{What I understand from this}

\begin{enumerate}
  \item The Stone construction gives the \emph{unconditional}
    object: $\hat{\mu}$ always exists.  It is the ``empirically
    adequate'' object in van Fraassen's sense.

  \item Descent to $\Omega$ is a \emph{choice}: which ultrafilters
    count as ``realized'' is additional structure over the algebra.

  \item The three-layer nesting:
    \[
      \pure(\Omega) \;\subseteq\;
      \{\text{countably complete ultrafilters}\} \;\subseteq\;
      \St(\mathcal{C}).
    \]
    $\sigma$-additivity gives mass on the middle set (intrinsic).
    Descent needs mass on the inner set (depends on $\Omega$).

  \item The van Fraassen identification (Stone measure $=$
    empirical adequacy (EA), descent $=$ realist commitment)
    has no precedent in the literature.  It is a philosophical
    contribution, not a mathematical one.
\end{enumerate}

% ==================================================================
\section{Layer 3: From probability to value-definiteness}
\label{sec:vdr}
% ==================================================================

\subsection{The classical results (all known)}

\begin{itemize}
  \item \textbf{Boolean algebras admit 2-valued homomorphisms:}
    Ultrafilters on Boolean algebras are exactly the 2-valued
    homomorphisms.  Existence by Zorn.
    \emph{Stone (1936); every Boolean algebra textbook.}

  \item \textbf{Kochen--Specker:} For the lattice $L(H)$ of closed
    linear subspaces of a Hilbert space $H$ with $\dim H \geq 3$,
    no dispersion-free state (2-valued homomorphism) exists
    \cite{KochenSpecker1967}.

  \item \textbf{Gleason:} For $L(H)$, $\dim H \geq 3$, every
    $\sigma$-additive state is $s(P) = \mathrm{tr}(\rho P)$
    \cite{Gleason1957}.

  \item \textbf{Bub--Clifton:} Given a state $\rho$ and a
    preferred observable $R$, there is a unique maximal Boolean
    subalgebra $D(\rho, R) \subseteq L(H)$ on which the state is
    a mixture of dispersion-free states.
    \cite{BubClifton1996, HalvorsonClifton1999}.

  \item \textbf{Orthomodular lattice (OML) duality:} The category
    of OMLs is dually equivalent to a category of orthomodular
    spaces.  Principal filters $P(A)$ are part of the dual
    structure (unlike Stone duality where they are invisible).
    \cite{McDonaldBimbo2023}.  Zero citations in quantum
    foundations as of 2026.

  \item \textbf{Colimit generation:} Gluing Boolean algebras via
    categorical pushout produces orthomodular lattices, generically
    non-distributive.
    \cite{Gunji2026}.
\end{itemize}

\subsection{What I proved (classical + one novel framing)}

\textbf{Commensurability theorem:} Distributivity is sufficient
for value-definite realism (VDR); for $L(H)$, $\dim H \geq 3$,
non-distributivity blocks it.  Lean-verified (0~sorry, 1~axiom for
Kochen--Specker (KS)).  \emph{The mathematics is the assembly of
Stone + KS + Gleason + Bub--Clifton.
Not a new theorem.}

\textbf{The EA / probabilistic realism (PR) / VDR vocabulary}
connecting van Fraassen's empirical adequacy to quantum logic via
McDonald--Bimb\'{o} duality is novel framing.  No precedent found.

\textbf{The coherence filtration}
$\mathcal{S}(A) \supseteq \mathcal{S}_\sigma(A) \supseteq
\mathcal{S}_{\mathrm{df}}(A)$ with three named transitions
(pasting, regularity, sharpness) organizes known results.  The
chain itself is textbook~(Pt\'{a}k--Pulmannov\'{a},
Kalmbach); the transition names and the observation that
distributivity collapses the filtration are new packaging.

\subsection{What I understand from this}

\begin{enumerate}
  \item \textbf{Distributivity controls value-definiteness.}
    Boolean algebras have ultrafilters (VDR available).
    Non-Boolean OMLs (dim $\geq 3$) do not (KS).  The dividing
    line is distributivity --- sufficient but not necessary in
    full generality (dim~2, horizontal sums).

  \item \textbf{The algebra can constrain realization in the OML
    case.} Unlike Boolean Stone duality, the McDonald--Bimb\'{o}
    dual carries $P(A)$ as structure.  ``Realization is additional
    structure'' in the Boolean case; it is \emph{constrained}
    structure in the OML case.

  \item \textbf{The realism debate has different mathematical
    answers for different algebras.}  Boolean: EA, PR, VDR all
    available; philosophical choice.  OML ($\dim \geq 3$): VDR
    blocked; PR available (Gleason); EA always available.

  \item \textbf{The topos programme is more powerful.}
    Isham--Butterfield / D\"{o}ring--Isham / Heunen--Landsman--Spitters
    encode the same structure via spectral presheaves and
    Bohrification.  Our approach is coarser (single dual space,
    not presheaf over contexts) but connects more directly to
    van Fraassen.
\end{enumerate}

% ==================================================================
\section{Delay reconstruction}
\label{sec:reconstruction}
% ==================================================================

\subsection{The classical results}

\begin{itemize}
  \item \textbf{Takens:} Generic smooth observable gives an
    injective delay map for $L \geq 2 \dim X$ \cite{Takens1981}.

  \item \textbf{Sauer--Yorke--Casdagli:} Prevalence-based
    embedding theorem \cite{SauerYorkeCasdagli1991}.

  \item \textbf{Generating partitions:}
    $h(T, P) = h(T)$ iff $P$ generates
    \cite{Sinai1959, Krieger1970}.

  \item \textbf{Grassberger--Procaccia:} Correlation integral
    estimates $D_2$ and $K_2$ (R\'{e}nyi-2 dimension and entropy)
    from delay data.

  \item \textbf{False nearest neighbours (FNN):} Embedding
    dimension selection \cite{KennelBrownAbarbanel1992}.  Sensitive
    to noise \cite{RhodesMorari1997}.

  \item \textbf{Ragwitz--Kantz:} Local prediction error
    minimization for $(d, L)$ selection \cite{RagwitzKantz2002}.

  \item \textbf{Botvinick-Greenhouse et al.:}
    Measure-theoretic Takens via optimal transport
    \cite{BotvinickGreenhouse2025}.
\end{itemize}

\subsection{What I proved (mostly classical)}

\textbf{Reconstruction theorem:} Observable $\sigma$-algebra $=$
full $\sigma$-algebra mod $\mu$ iff delay algebra dense in $L^2$
iff delay map is measure-theoretic embedding.  This is
Krieger's generating partition theorem~\cite{Krieger1970} in
delay-map language.

\textbf{Predictive factorization, Chapman--Kolmogorov, Koopman--Perron:}
All classical (Rokhlin, Kolmogorov, Koopman 1931).

\textbf{Finite-sample convergence rates:} Minimax rates
$n^{-s/(2s+d)}$~\cite{Stone1982}; elbow stopping via
Lepski~\cite{Lepski1991}.  Not cited in the original work ---
a significant attribution gap.

\subsection{What is genuinely novel (but motivation unclear)}

\textbf{Geometric $\neq$ algebraic reconstruction (2026-05-14):}
The bridge note claimed $\delta(L) \to 0 \Leftrightarrow
(\mu \otimes \mu)(R_L) \to 0$ under a fibre mixing condition.
This is \textbf{false}: the core identity (Lemma~2.1 of the
bridge note) has a disintegration error.  The conditional
variance $\mathbb{E}[\mathrm{Var}(\mathbf{1}_S \mid \mathcal{O}_L)]
= \sum_z p_z \,\mu_z(S)\mu_z(S^c)$ is weighted by $p_z$,
while the collision integral $(1/2)\int_{R_L}
|\mathbf{1}_S - \mathbf{1}_{S'}|^2\,d(\mu\otimes\mu)
= \sum_z p_z^2 \,\mu_z(S)\mu_z(S^c)$ is weighted by $p_z^2$.
The bridge note equated them; they are genuinely different.

\textbf{Skew-product counterexample:} $X = A^{\mathbb{Z}} \times
B^{\mathbb{Z}}$, $T = \sigma \times \sigma$, $h(a,b) = a_0$,
independent iid uniform.  Then collision $= 2^{-(L+1)} \to 0$
(fibres shrink in the $a$-coordinate) but $\delta_L = 1/4$ for
all $L$ (the event $\{b_0 = 0\}$ is maximally unresolved).
This falsifies even the ``easy direction.''

\textbf{What survives:} The definitions (collision probability,
Rokhlin distance, fibre mixing) are well-posed.
The trivial direction $\delta \to 0 \Rightarrow
\text{collision} \to 0$ holds.  The converse and the entropy
characterization $\delta \to 0 \Leftrightarrow H_2 \to \infty$
are both dead.

\textbf{Replacement direction:} A short note proving that
geometric and algebraic reconstruction are inequivalent,
with the skew-product as the central theorem.  The open
question becomes: what conditions on the factor map make
collision $\to 0$ imply $\delta \to 0$?

\subsection{What I understand from this}

\begin{enumerate}
  \item The reconstruction problem is \emph{algebraically
    underdetermined}: the observation algebra does not constrain
    the realization.  Embedding selection ($d$, $L$) is necessarily
    empirical.

  \item Every practical reconstruction criterion (FNN, mutual
    information, Lyapunov convergence, prediction error) implicitly
    assumes a modelling commitment.  Deterministic criteria fail under noise because
    noise violates the commitment.

  \item The Rokhlin distance $\delta(\mathcal{O})$ is the
    measure-theoretic notion of ``how close is the observable
    $\sigma$-algebra to generating everything.''  It is the
    Rokhlin metric, named and studied since 1967.

  \item Collision probability / R\'{e}nyi-2 entropy measures
    the ``collision structure'' of the fibre decomposition.
    Grassberger--Procaccia already extracts this from data.

  \item The bridge theorem linking collision convergence to
    $\sigma$-algebra generation is false: the disintegration error
    ($p_z$ vs $p_z^2$ weighting) means the proposed identity is
    simply the wrong identity, and the inequivalence of geometric
    and algebraic reconstruction is known/obvious.
\end{enumerate}

% ==================================================================
\section{Dead ends and lessons}
\label{sec:dead-ends}
% ==================================================================

\begin{enumerate}
  \item \textbf{De Finetti via Stone:} Exchangeability forces
    $\sigma$-additivity (Hausdorff moments).  The ``exotic part''
    of the Stone measure is empty.  Dead end.

  \item \textbf{Stationary Stone:} The Yosida--Hewitt decomposition
    commutes with shift-invariance.  An observation, not a theorem.
    The real open question (ergodic theory of the purely finitely
    additive part) is out of reach.

  \item \textbf{Dynamics from observation:} ``Can the Stone space
    give rise to dynamics without positing $T$?''  No: the algebra
    doesn't constrain $\Omega$, so it certainly doesn't constrain
    structure on $\Omega$.

  \item \textbf{Paper III (noise thresholds):} Residual
    autocorrelation as embedding diagnostic =
    Billings--Voon~\cite{BillingsVoon1986} /
    Casdagli~\cite{Casdagli1992}.  FNN noise sensitivity =
    Rhodes--Morari~\cite{RhodesMorari1997}.  Elbow stopping =
    Lepski~\cite{Lepski1991}.  Rediscovery throughout.

  \item \textbf{``No intermediate condition'':} The claim that
    no condition interpolates between first-order and
    $\sigma$-additivity is false.
    $\limsup \ell(A_n) < c$ trivially interpolates.

  \item \textbf{Bridge theorem (fibre mixing):} The claimed
    equivalence $\delta(L) \to 0 \Leftrightarrow
    (\mu \otimes \mu)(R_L) \to 0$ under fibre mixing is false.
    The core identity has a disintegration error: the conditional
    variance integral is weighted by $p_z$, the collision integral
    by $p_z^2$.  The skew-product $X = A^{\mathbb{Z}} \times
    B^{\mathbb{Z}}$, $h(a,b) = a_0$ falsifies even the easy
    direction.  The replacement observation (geometric $\neq$
    algebraic reconstruction) is known to ergodic theorists.
    Both the positive bridge and the negative note fail the audit.

  \item \textbf{Entropy characterization:}
    $\delta(L) \to 0 \Leftrightarrow H_2(\nu_L) \to \infty$
    depended entirely on the bridge theorem.  Dead.

  \item \textbf{Nonlinear semigroups / Hille operator (Phase 1):}
    $C_n$-semigroups on Banach spaces, resolvent from finite
    difference spectra, max-entropy fitting.  Abandoned approach
    --- the semigroup formalism was too heavy for the applied
    question.  The Rose kernel was the lighter successor.

  \item \textbf{Rose--Takens measure conjugacy:}
    ``Takens embedding + $D_{\mathrm{KL}}(h_\#\mu \| \nu) = 0$
    implies measure conjugacy.''  Trivially true by definition:
    diffeomorphism + pushforward = target measure IS measure
    conjugacy.  Not a theorem.

  \item \textbf{CGF $\to$ Fisher $\to$ Landau chain:}
    Standard information geometry~\cite{Amari1985}.  The CGF-Fisher
    connection is the definition of Fisher information.

  \item \textbf{KL = least action (discrete Onsager--Machlup):}
    Known: \cite{Eyink1996}, \cite{MaesNetocny2003}, large
    deviations literature.

  \item \textbf{Fisher clock / entropy production:}
    Broken detailed balance $\to$ entropy production is
    \cite{Schnakenberg1976}.  The ``Fisher clock'' framing is
    poetic but the mathematics is standard nonequilibrium
    statistical mechanics.

  \item \textbf{Neutral operator edge law:}
    The joint edge law $\Gamma = \nu \otimes K$ and its
    $f$-divergence structure were absorbed into Paper II
    (conditional regularity kernel $\kappa_Q$).  The diagnostic
    exponents $\eta = 2s$, $\zeta = 2s/(2s+d)$ are the same
    Lepski/Stone rates already canned in Paper III.
\end{enumerate}

% ==================================================================
\section{What remains genuinely open}
\label{sec:open}
% ==================================================================

\begin{enumerate}
  \item \textbf{OML extension problem:} Does every state on a
    general OML extend to a $\sigma$-additive measure on the
    McDonald--Bimb\'{o} dual $S_0(A)$?  Gleason answers for
    $L(H)$; the general case is \textbf{no} --- Pt\'{a}k (1987)
    ``Exotic logics'' constructs $\sigma$-orthocomplete OMPs
    admitting finitely additive states but no $\sigma$-additive
    states.  The structural mechanism: the center $C(L)$ can be
    a Boolean algebra with no $\sigma$-additive probability,
    blocking extension.

    \emph{Remaining open question:} what condition on a non-Boolean
    OML forces $\sigma$-additivity?  For $L(H)$ ($\dim \geq 3$),
    the center is trivial and Gleason provides it from lattice
    geometry.  For general OMLs with trivial center but discrete
    structure, the question is open.

  \item \textbf{Ergodic theory of purely finitely additive charges:}
    What are the extreme points of the purely finitely additive
    shift-invariant charges?  Connected to Banach limits.
    \emph{Out of reach currently.}
\end{enumerate}

% ==================================================================
\section{The understanding, condensed}
\label{sec:condensed}
% ==================================================================

One sentence per insight:

\begin{enumerate}
  \item The passage from coherent charges to probability requires
    $\sigma$-additivity, which is not derivable from finitary
    conditions.

  \item The Stone construction gives an unconditional measure on
    the dual space; descent to a realization requires a commitment
    the algebra cannot force.

  \item Which ultrafilters are ``realized'' is additional structure
    in the Boolean case and constrained structure in the OML case.

  \item Distributivity is sufficient for value-definiteness;
    non-distributivity blocks it for $L(H)$, $\dim \geq 3$, but
    not for all non-distributive OMLs (e.g.\ $L(\mathbb{C}^2)$,
    horizontal sums).

  \item The realism/empiricism boundary is algebraic:
    van Fraassen's empirical adequacy is the unconditional level;
    value-definite realism requires the algebra to admit
    dispersion-free states (guaranteed by distributivity,
    blocked by KS for $L(H)$, $\dim \geq 3$).

  \item Delay reconstruction is algebraically underdetermined:
    embedding selection is necessarily empirical.

  \item Geometric reconstruction (collision $\to 0$) and algebraic
    reconstruction ($\delta \to 0$) are genuinely inequivalent.
    The bridge theorem claiming their equivalence under fibre
    mixing was false (disintegration error in the core identity).
    The skew-product counterexample is definitive.

  \item Every practical reconstruction diagnostic implicitly
    assumes a modelling commitment; noise violates deterministic
    commitments.

  \item Audit before you draft.
\end{enumerate}

% ==================================================================
\section{References by subject}
\label{sec:refs}
% ==================================================================

\subsection*{Measure extension and charges}
\begin{itemize}
  \item Halmos, \emph{Measure Theory}, Van Nostrand, 1950.  \S13 (extension), \S53--54 (Boolean algebras).~\cite{Halmos1950}
  \item Billingsley, \emph{Probability and Measure}, 3rd ed., Wiley, 1995.  Ch.~36 (Kolmogorov extension).~\cite{Billingsley1995}
  \item Bogachev, \emph{Measure Theory}, Vols.~1--2, Springer, 2007.  Vol.~1 Ch.~1 (Carath\'{e}odory).~\cite{Bogachev2007}
  \item Rao \& Rao, \emph{Theory of Charges}, Academic Press, 1983.  Thm~10.5.3 (Stone-space support characterization).~\cite{RaoRao1983}
  \item Yosida \& Hewitt, ``Finitely additive measures,'' \emph{Trans.\ AMS} 72 (1952), 46--66.~\cite{YosidaHewitt1952}
  \item Choksi, ``Inverse limits of measure spaces,'' \emph{Proc.\ London Math.\ Soc.} s3-8 (1958), 321--342.  Thm~1.~\cite{Choksi1958}
  \item Fremlin, \emph{Measure Theory}, Vol.~3, Torres Fremlin, 2003.  \S326 (additive functionals on Boolean algebras).~\cite{Fremlin2003}
\end{itemize}

\subsection*{Stone duality}
\begin{itemize}
  \item Johnstone, \emph{Stone Spaces}, Cambridge University Press, 1982.~\cite{Johnstone1982}
\end{itemize}

\subsection*{Quantum logic and contextuality}
\begin{itemize}
  \item Kalmbach, \emph{Orthomodular Lattices}, Academic Press, 1983.  Ch.~4 (states, valuations).~\cite{Kalmbach1983}
  \item Pt\'{a}k \& Pulmannov\'{a}, \emph{Orthomodular Structures as Quantum Logics}, Kluwer, 1991.~\cite{PtakPulmannova1991}
  \item Kochen \& Specker, ``The problem of hidden variables in quantum mechanics,'' \emph{J.\ Math.\ Mech.} 17 (1967), 59--87.~\cite{KochenSpecker1967}
  \item Gleason, ``Measures on the closed subspaces of a Hilbert space,'' \emph{J.\ Math.\ Mech.} 6 (1957), 885--893.~\cite{Gleason1957}
  \item Bub \& Clifton, ``A uniqueness theorem for `no collapse' interpretations,'' \emph{Stud.\ Hist.\ Phil.\ Mod.\ Phys.} 27 (1996), 181--219.~\cite{BubClifton1996}
  \item Halvorson \& Clifton, ``Maximal beable subalgebras,'' \emph{Int.\ J.\ Theor.\ Phys.} 38 (1999), 2441--2484.~\cite{HalvorsonClifton1999}
  \item Abramsky \& Brandenburger, ``The sheaf-theoretic structure of non-locality and contextuality,'' \emph{New J.\ Phys.} 13 (2011), 113036.~\cite{AbramskyBrandenburger2011}
  \item McDonald \& Bimb\'{o}, ``Topological duality for orthomodular lattices,'' \emph{Math.\ Logic Quart.} 69 (2023), 514--541.~\cite{McDonaldBimbo2023}
  \item Gunji, Haruna \& Uragami, ``Contextuality as a left adjoint,'' arXiv:2603.22353 (2026).~\cite{Gunji2026}
\end{itemize}

\subsection*{Topos approaches}
\begin{itemize}
  \item Isham \& Butterfield, ``Topos perspective on the Kochen--Specker theorem: I,'' \emph{Int.\ J.\ Theor.\ Phys.} 37 (1998), 2669--2733.~\cite{IshamButterfield1998}
  \item Heunen, Landsman \& Spitters, ``A topos for algebraic quantum theory,'' \emph{Comm.\ Math.\ Phys.} 291 (2009), 63--110.~\cite{HeunenLandsmanSpitters2009}
\end{itemize}

\subsection*{Delay embedding and reconstruction}
\begin{itemize}
  \item Takens, ``Detecting strange attractors in turbulence,'' \emph{Lecture Notes in Math.} 898 (1981), 366--381.~\cite{Takens1981}
  \item Sauer, Yorke \& Casdagli, ``Embedology,'' \emph{J.\ Stat.\ Phys.} 65 (1991), 579--616.~\cite{SauerYorkeCasdagli1991}
  \item Kennel, Brown \& Abarbanel, ``Determining embedding dimension for phase-space reconstruction,'' \emph{Phys.\ Rev.\ A} 45 (1992), 3403--3411.~\cite{KennelBrownAbarbanel1992}
  \item Ragwitz \& Kantz, ``Markov models from data by simple nonlinear time series predictors,'' \emph{Phys.\ Rev.\ E} 65 (2002), 056201.~\cite{RagwitzKantz2002}
  \item Botvinick-Greenhouse, Oprea \& Maulik, ``Measure-theoretic time-delay embedding,'' \emph{J.\ Stat.\ Phys.} 192 (2025), 171.~\cite{BotvinickGreenhouse2025}
  \item Rhodes \& Morari, ``False-nearest-neighbors algorithm and noise-corrupted time series,'' \emph{Phys.\ Rev.\ E} 55 (1997), 6162--6170.~\cite{RhodesMorari1997}
\end{itemize}

\subsection*{Ergodic theory}
\begin{itemize}
  \item Rokhlin, ``Lectures on the entropy theory of measure-preserving transformations,'' \emph{Russian Math.\ Surveys} 22 (1967), 1--52.~\cite{Rokhlin1967}
  \item Sinai, ``On the concept of entropy of a dynamical system,'' \emph{Dokl.\ Akad.\ Nauk SSSR} 124 (1959), 768--771.~\cite{Sinai1959}
  \item Krieger, ``On entropy and generators of measure-preserving transformations,'' \emph{Trans.\ AMS} 149 (1970), 453--464.~\cite{Krieger1970}
  \item Furstenberg, ``Ergodic behavior of diagonal measures and a theorem of Szemer\'{e}di,'' \emph{J.\ Analyse Math.} 31 (1977), 204--256.~\cite{Furstenberg1977}
\end{itemize}

\subsection*{Philosophy of science}
\begin{itemize}
  \item van Fraassen, \emph{The Scientific Image}, Oxford University Press, 1980.~\cite{vanFraassen1980}
  \item de Finetti, \emph{Theory of Probability}, Vol.~1, Wiley, 1974.~\cite{deFinetti1974}
\end{itemize}

\subsection*{Nonparametric statistics}
\begin{itemize}
  \item Stone, ``Optimal global rates of convergence for nonparametric regression,'' \emph{Ann.\ Statist.} 10 (1982), 1040--1053.~\cite{Stone1982}
  \item Lepski, ``On a problem of adaptive estimation in Gaussian white noise,'' \emph{Theory Probab.\ Appl.} 35 (1991), 454--466.~\cite{Lepski1991}
  \item Billings \& Voon, ``Correlation based model validity tests for non-linear models,'' \emph{Int.\ J.\ Control} 44 (1986), 235--244.~\cite{BillingsVoon1986}
  \item Casdagli, ``Chaos and deterministic versus stochastic non-linear modelling,'' \emph{J.\ Roy.\ Statist.\ Soc.\ B} 54 (1992), 303--328.~\cite{Casdagli1992}
\end{itemize}

\subsection*{Model theory and probability logic}
\begin{itemize}
  \item Hoover, ``Probability logic,'' \emph{Ann.\ Math.\ Logic} 14 (1978), 287--313.~\cite{Hoover1978}
  \item Fagin, Halpern \& Megiddo, ``A logic for reasoning about probabilities,'' \emph{Inform.\ Comput.} 87 (1990), 78--128.~\cite{FaginHalpernMegiddo1990}
\end{itemize}

